\section*{Chapitre \ref{ch:b3:total-synthesis} --- Synthèse totale~: stratégie et voies classiques}

\begin{solution}{exo:b3:total-synthesis:1}
À 90\,\%~: $0.9^{10} = 35$\,\% et $0.9^{20} = 12$\,\%. À 80\,\%~: $0.8^{10} = 11$\,\% et $0.8^{20} = 1.2$\,\%.
\end{solution}

\begin{solution}{exo:b3:total-synthesis:2}
$5 + 3 + 1 + 4 = 13$ étapes en tout~; LLS $5 + 1 + 4 = 10$.
\end{solution}

\begin{solution}{exo:b3:total-synthesis:3}
4-(Diméthylamino)butan-2-one, $\ce{CH3COCH2CH2N(CH3)2}$~: l'énol de l'acétone s'additionne sur l'ion
iminium $\ce{CH2=N+(CH3)2}$.
\end{solution}

\begin{solution}{exo:b3:total-synthesis:4}
Déconnecter les deux liaisons du cycle en position allylique par rapport à
la liaison C=C~: le buta-1,3-diène et la but-3-én-2-one (méthylvinylcétone),
qui réagissent par chauffage.
\end{solution}

\begin{solution}{exo:b3:total-synthesis:5}
Linéaire~: $0.9^{21} = 11$\,\%. Convergente~: $0.9^{11} = 31$\,\%, presque trois fois plus.
\end{solution}

\begin{solution}{exo:b3:total-synthesis:6}
A~: $(7 + 1)/12 = 67$\,\%. B~: $7/9 = 78$\,\%.
\end{solution}

\begin{solution}{exo:b3:total-synthesis:7}
Sans~: $0.9^8 = 43$\,\%. Avec quatre étapes supplémentaires à 95\,\%~: $43\,\% \times 0.95^4 = 35$\,\%.
\end{solution}

\begin{solution}{exo:b3:total-synthesis:8}
Primaire~: un éther silylé encombré (\textit{tert}-butyldiphénylsilyle),
retiré par les ions fluorure. Secondaire~: un éther de 4-méthoxybenzyle,
retiré par oxydation au DDQ. L'alcool tertiaire, encombré, peut souvent
rester libre, ou être protégé sous forme d'un petit éther silylé coupé par
un acide doux. Chaque condition laisse intacts les autres groupes.
\end{solution}

\begin{solution}{exo:b3:total-synthesis:9}
Les \omterm{def:b3:total-synthesis:total}{synthèses totales} demandaient des dizaines d'étapes avec des
rendements globaux bien inférieurs à 1\,\%~; un composé possédant tout le
système de cycles et la plupart des stéréocentres, extrait des aiguilles
de l'if commun européen, source renouvelable, est converti en quelques
étapes.
\end{solution}

\begin{solution}{exo:b3:total-synthesis:10}
Addition de Michael~: l'énolate de la dione (C2, entre les carbonyles)
s'additionne sur le carbone $\beta$ de la méthylvinylcétone, ce qui donne une tricétone. Aldolisation~: l'énolate de la
méthylcétone attaque un carbonyle du cycle de façon intramoléculaire, et la
déshydratation donne la cyclohexénone~: la cétone de Wieland--Miescher, une
ènedione bicyclique portant un groupe méthyle angulaire.
\end{solution}

\begin{solution}{exo:b3:total-synthesis:11}
$y^{m+1}/y^{2m+1} = y^{-m}$~: à 90\,\%, 1.7 pour $m = 5$ et 4.9 pour $m = 15$. Plus la voie est longue, plus la
convergence paie.
\end{solution}

\begin{solution}{exo:b3:total-synthesis:12}
Chaque cycle se ferme avec le cation et la double liaison attaquante dans
une disposition de type chaise~; l'addition sur chaque double liaison est
anti, et les substituants prennent des positions pseudo-équatoriales, ce
qui rend \textit{trans} chaque jonction de cycles. Une double liaison \textit{E}
place ses deux prolongements de chaîne de sorte que la jonction \textit{trans} en résulte~; une double liaison \textit{Z}
donnerait la jonction \textit{cis}~: la géométrie du polyène code la
stéréochimie du produit.
\end{solution}

\begin{solution}{pb:b3:total-synthesis:1}
\textbf{1.} $\ce{C4H6O2}$, $\ce{CH5N}$, $\ce{C5H6O5}$, $\ce{C8H13NO}$.

\textbf{2.} $\ce{C4H6O2 + CH5N + C5H6O5 -> C8H13NO + 2CO2 + 2H2O}$.

\textbf{3.} 86.0, 31.0 et \qty{146.0}{g/mol}~; tropinone \qty{139.0}{g/mol}.

\textbf{4.} $139.0/(86.0 + 31.0 + 146.0) = 0.53$~: 53\,\%.

\textbf{5.} Ses groupes $\ce{CH2}$ centraux, situés chacun entre deux groupes carbonyle, s'énolisent facilement dans
l'eau~; l'acétone elle-même est trop faiblement nucléophile, et les groupes
carboxyle activants partent ensuite.

\textbf{6.} La formation des ions iminium exige de l'amine libre (trop
d'acide la protone), et l'énol comme la formation de l'iminium exigent un
peu d'acide~: un pH proche de la neutralité les concilie.

\textbf{7.} La méthylamine, $\ce{CH3NH2}$, s'additionne sur l'un des groupes aldéhyde~; la perte d'une molécule d'eau donne ensuite l'ion iminium
$\ce{R-CH=N+H-CH3}$.

\textbf{8.} Un carbone d'énol de l'acide acétonedicarboxylique attaque le
carbone de l'iminium, ce qui forme une liaison C--C et une amine
secondaire.

\textbf{9.} L'amine secondaire se condense avec le second groupe aldéhyde de
la même chaîne, ce qui donne un ion iminium cyclique (à cinq chaînons).

\textbf{10.} L'énol situé de l'autre côté de la cétone l'attaque, ce qui
ferme le cycle à six chaînons et le squelette bicyclique.

\textbf{11.} Chacun est un $\beta$-céto-acide, qui perd $\ce{CO2}$ via un état de transition cyclique à six
chaînons, ce qui donne l'énol de la cétone.

\textbf{12.} Deux liaisons C--C et deux liaisons C--N.

\textbf{13.} $0.75^{14} = 0.018$~: 1.8\,\%.

\textbf{14.} $42/1.78 = 24$.

\textbf{15.} 100\,\%~: une seule étape, entièrement consacrée à la construction.

\textbf{16.} Une seule opération, l'eau comme solvant, des composants bon
marché, aucun groupe protecteur et presque aucun déchet hormis le dioxyde
de carbone et l'eau.

\textbf{17.} La plante construit le squelette tropane à partir d'un ion
iminium cyclique dérivé d'un acide aminé et d'une unité dérivée de
l'acétate, par le même type de chimie de Mannich.

\textbf{18.} Les cascades (et les \omterm{def:b3:pericyclic:families}{cycloadditions} formatrices de cycles comme
la réaction de Diels--Alder).

\textbf{19.} Non~: un plan miroir passe par N, son groupe méthyle, C3 et
l'oxygène du carbonyle.

\textbf{20.} Chacun porte quatre groupes différents, mais ils sont images
l'un de l'autre par le \omterm{def:b3:point-groups:mirror}{plan de symétrie} propre de la molécule~: un composé
de type méso.

\textbf{21.} Des diastéréoisomères, qui diffèrent par l'orientation du OH
porté par C3 par rapport au pont azoté.

\textbf{22.} Non~: le plan miroir est conservé dans les deux.

\textbf{23.} Les deux faces du groupe carbonyle ne sont pas équivalentes~:
l'une fait face au pont azoté, l'autre au pont à deux carbones, qui gênent
différemment.

\textbf{24.} La voie monotope donne un \omterm{def:b3:total-synthesis:architecture}{rendement global} environ 24 fois plus
élevé que la voie linéaire de 14 étapes.
\end{solution}
